\documentclass{tufte-handout}

\usepackage[T1]{fontenc}
\usepackage[utf8]{inputenc}
\usepackage{graphicx}
\usepackage{booktabs}
\usepackage{natbib}
\usepackage{hyperref}
\usepackage{listings}
\usepackage{color}

\graphicspath{{/home/mwl04747/RTricks/22_Biogeochemistry_Template/figures/}}

\lstset{
  language=R,
  basicstyle=\small\ttfamily,
  commentstyle=\color{gray},
  breaklines=true
}

\bibliographystyle{plainnat}

% ── Fill in your details ──────────────────────────────────────────────────────
\title{Report Title}
\author{Your Name}
\date{\today}

% ══════════════════════════════════════════════════════════════════════════════
\begin{document}
\maketitle

\begin{abstract}
\noindent One-paragraph summary of your topic, methods, and findings \citep{R-base}.
\end{abstract}

% ──────────────────────────────────────────────────────────────────────────────
\section{Introduction}

\newthought{Write your introduction here.}
Use sidenotes for supplementary details.\sidenote{Like this. Keep them brief.}

% ──────────────────────────────────────────────────────────────────────────────
\section{Graphics}

% Full-width figure — spans text column and margin
\begin{figure*}
  \includegraphics[width=\linewidth]{your-plot}% filename in figures/
  \caption{Full-width figure caption. Describe the key finding.}
\end{figure*}

% Margin figure — sits in the right margin
\begin{marginfigure}
  \includegraphics[width=\linewidth]{your-margin-plot}
  \caption{Margin figure caption.}
\end{marginfigure}

% Body-width figure
\begin{figure}
  \includegraphics[width=\linewidth]{your-body-plot}
  \caption{Body-width figure caption.}
\end{figure}

Save R plots to the \texttt{figures/} folder like so:

\begin{lstlisting}
library(ggplot2)
p <- ggplot(mtcars, aes(wt, mpg)) + geom_point() +
       labs(x = "Weight (1000 lbs)", y = "MPG")
ggsave("figures/your-plot.pdf", p, width = 6, height = 4)
\end{lstlisting}

% ──────────────────────────────────────────────────────────────────────────────
\section{Analysis}

Write your analysis here \citep{rmarkdown2018, rmarkdown2020}.
Inline math: $\bar{x} = \frac{1}{n}\sum x_i$. Display math:
\[
  s^{2} = \frac{1}{n-1}\sum_{i=1}^{n}(x_i - \bar{x})^{2}
\]

\begin{table}
  \centering
  \caption{Summary table --- replace with your data.}
  \begin{tabular}{lrrr}
    \toprule
    Group & $n$ & Mean & SD \\
    \midrule
    A & 50 & 12.3 & 2.1 \\
    B & 48 & 14.7 & 1.9 \\
    \bottomrule
  \end{tabular}
\end{table}

% ──────────────────────────────────────────────────────────────────────────────
\section{Conclusion}

Summarise key findings here.

% ──────────────────────────────────────────────────────────────────────────────
\section{Glossary}

\begin{description}
  \item[\textbf{R}] A free, open-source language and environment for statistical
    computing and graphics, maintained by the R Foundation \citep{R-base}.
  \item[\textbf{R Markdown}] A file format (\texttt{.Rmd}) that combines prose,
    R code, and output into a single reproducible document \citep{R-rmarkdown}.
  \item[\textbf{Tufte Handout}] A document class inspired by Edward Tufte's
    books, featuring wide margins for sidenotes and figures.
  \item[\textbf{Sidenote}] A marginal annotation that supplements body text
    without interrupting reading flow; created with \texttt{\textbackslash sidenote\{\}}.
  \item[\textbf{Reproducible Research}] The practice of embedding code, data,
    and narrative together so results can be independently verified and rerun.
  \item[\textbf{\texttt{graphicspath}}] A \LaTeX{} command that tells the
    compiler where to look for image files, avoiding full paths in every
    \texttt{\textbackslash includegraphics} call.
  % ── Add your own terms below ─────────────────────────────────────────────
  \item[\textbf{Term 7}] Your definition here.
\end{description}

% ──────────────────────────────────────────────────────────────────────────────
\bibliography{references}

\end{document}